% Part: turing-machines
% Chapter: undecidability
% Section: universal-tm

\documentclass[../../../include/open-logic-section]{subfiles}

\begin{document}

\olfileid{tur}{und}{uni}
\olsection{ہمہ گیر ٹیورنگ مشیناں}

\olref[enu]{sec} وچ اسی گل کیتی سی کہ ہر ٹیورنگ مشین دا
بیان صحیح اعداد دے اک محدود سلسلے نال دتا جا سکدا اے۔
ایہ سلسلہ مشین دیاں حالتاں، حرفی نشاناں، ابتدائی حالت
تے ہدایات نوں کوڈ کردا اے۔ اسی ایہ وی دَسیا سی کہ
ایہناں سارے بیانواں دا مجموعہ !!{enumerable} اے۔ کیوں
جے ایہناں بیانواں دا مجموعہ !!{denumerable} اے، ایس دا
مطلب اے کہ~$\Nat$ توں ایہناں بیانواں تک اک
!!a{surjective} فنکشن موجود اے۔ ایہو جیہا !!a{surjective}
فنکشن، مثلاً، کینٹر دے ٹیڑھے-میڑھے طریقے نال لبھ
سکدا اے۔ ایہ سانوں ساریاں ٹیورنگ مشیناں (دے بیانواں)
دی گنتی وار فہرست بناؤن دا طریقہ دیندا اے۔ ایہو جیہی اک
فہرست مقرر کر لئیے، تاں $1$ نمبر، $2$ نمبر، \dots، $e$ نمبر
والی ٹیورنگ مشین دی گل کر سکدے آں۔ ایہناں
عدداں نوں \emph{اشاریے} کہندے نیں۔

\begin{defn}
جے $M$~ساڈی مقررہ فہرست وچ $e$ نمبر والی ٹیورنگ مشین اے،
تاں اسی $e$~نوں~$M$ دا اک \emph{اشاریہ} کہندے آں۔
اسی $M_e$ لکھ کے $e$ نمبر والی ٹیورنگ مشین مراد لیندے آں۔
\end{defn}

اک مشین دے کئی اشاریے ہو سکدے نیں۔ مثلاً~$M$ دے
دو بیانواں وچ اوہدیاں ہدایات فہرست وچ لکھن دی ترتیب
وکھ ہو سکدی اے، تے ایہناں وکھرے بیانواں دے اشاریے
وی وکھ ہون گے۔

اہم گل ایہ اے کہ ٹیورنگ مشیناں دے بیانواں دی فہرست
ایس طرح بنائی جا سکدی اے کہ اشاریے توں~$M$ دا بیان
مؤثر طریقے نال حساب ہو جاوے، تے مشین~$M$ دے بیان توں
اوس دا اک اشاریہ وی مؤثر طریقے نال حساب ہو جاوے۔
فیر چرچ--ٹیورنگ دے دعوے مطابق ایسی ٹیورنگ مشین
لبھ سکدی اے جیڑی اشاریے~$e$ والی مشین دا بیان
حاصل کرے تے اوہ بیان اپنی پٹی اُتے آؤٹ پٹ دے طور
اُتے لکھ دے۔ ایہ بیان~$\TMstroke$ دیاں لڑیوں
دے سلسلے دی صورت وچ ہووے گا (ہر لڑی~$M_e$
دے بیان والے سلسلے دے اک مثبت صحیح عدد دی نمائندگی کرے گی)۔

ایس توں بعد ایہ سوال فطری اے: ٹیورنگ مشین دے اشاریے
اُتے تعریف ہون والے کیہڑے فنکشن خود ٹیورنگ مشین نال
حساب ہو سکدے نیں؟ ٹیورنگ مشین دے اشارییاں دیاں
کیہڑیاں خاصیتاں دا فیصلہ ٹیورنگ مشین کر سکدی اے؟
اک مثال: جیہڑا فنکشن اشاریے~$e$ نوں اوس مشین دیاں
حالتاں دی تعداد ول بھیجدا اے جیہدا اشاریہ~$e$ اے،
اوہ ٹیورنگ مشین نال حساب ہو سکدا اے۔ ایسی مشین ایہ کم
کرے گی: جے اوہنوں پٹی اُتے $e$~$\TMstroke$ نشاناں دی
اک لڑی دتی جاوے، تاں پہلے اوہ $e$ دا کوڈ کھول کے
مشین دا بیان حاصل کرے گی۔ ہن ایہ بیان پٹی اُتے
~$\TMstroke$ دیاں لڑیوں دے سلسلے نال ظاہر اے۔ ایس
سلسلے دا پہلا !!{element} حالتاں دی تعداد اے۔ ایس لئی
ہن مشین نوں پہلی $\TMstroke$ لڑی توں علاوہ باقی سب
مٹانا تے فیر رُکنا اے۔

اک قابلِ توجہ نتیجہ ایہ اے:

\begin{thm}\ollabel{thm:universal-tm} اک \emph{ہمہ گیر
  ٹیورنگ مشین}~$U$ موجود اے جیڑی انپٹ $\tuple{e,n}$
  نال شروع ہووے تاں
  \begin{enumerate}
    \item اُدوں تے صرف اُدوں رُکدی اے جدوں $M_e$ انپٹ~$n$
    اُتے رُکدی اے، تے
    \item جے $M_e$ آؤٹ پٹ $m$ نال رُکے، تاں~$U$ وی
    اوہو آؤٹ پٹ لے کے رُکدی اے۔
  \end{enumerate}
  ایس طرح $U$ اوہ فنکشن $f\colon \Nat \times \Nat \pto \Nat$
  دا حساب کردی اے جیہدے لئی $f(e,n) = m$ ہووے جے $M_e$
  انپٹ~$n$ نال شروع ہو کے آؤٹ پٹ~$m$ دے کے رُکے،
  تے دوجی صورت وچ ایہ تعریف شدہ نہ ہووے۔
\end{thm}

\begin{proof}
  $U$ نوں حقیقت وچ مکمل بنا کے پیش کرنا عملاً ناممکن
  جیہا اے، کیوں جے ایہ بے حد پیچیدہ مشین اے۔ پر اسی
  اوہدے عمل دا خاکہ بیان کر سکدے آں تے فیر چرچ--ٹیورنگ
  دے دعوے نوں لاگو کر سکدے آں۔ شروع وچ $U$ دی پٹی
  اُتے $e$ عدد وچ $\TMstroke$ نشاناں دی اک لڑی اے،
  تے اوس دے بعد $n$~$\TMstroke$ نشاناں دی لڑی اے۔
  مشین پہلے اشاریے~$e$ دا کوڈ انپٹ~$n$ دے سجّے
  پاسے کھولدی اے۔ ایس نال اعداد دی اک فہرست ملدی
  اے، یعنی~$\TMstroke$ دیاں لڑیوں جیڑیاں~$\TMblank$
  نال وکھ نیں، تے ایہ مشین~$M_e$ دیاں ہدایات بیان
  کردیاں نیں۔ فیر $U$، مشین~$M_e$ دی ابتدائی حالت
  دا عدد تے عدد~$1$، مشین~$M_e$ دے بیان دے سجّے
  پاسے لکھدی اے۔ (ایہ وی اک نشان والی نمائندگی وچ
  $\TMstroke$ دیاں لڑیوں نیں۔) اگّے اوہ انپٹ،
  یعنی $n$~$\TMstroke$ نشاناں والی لڑی، سجّے
  پاسے نقل کردی اے؛ پر ہر $\TMstroke$ دی تھاں
  تین $\TMstroke$ نشاناں دی لڑی رکھدی اے
  (چیتے رکھو کہ $\TMstroke$ نشان دا عدد~$3$ اے؛
  $1$،~$\TMendtape$ دا عدد، تے $2$،~$\TMblank$
  دا عدد اے)۔ ایہناں لڑیوں دے سلسلے دے کھبے
  سِرے اُتے، جیڑیاں~$\TMblank$ نال $U$ دی پٹی
  اُتے وکھ نیں، اوہ اکّو~$\TMstroke$ لکھدی اے؛
  ایہ~$\TMendtape$ دا کوڈ اے۔

  ہن $U$ دی پٹی اُتے ایہ سب کجھ اے: اشاریہ~$e$،
  عدد~$n$، ابتدائی حالت دا کوڈ عدد (یعنی “موجودہ
  حالت”)، سر دی ابتدائی جگہ~$1$ دا عدد (یعنی
  “سر دی موجودہ جگہ”)، تے مصنوعی “پٹی” دا ابتدائی
  مواد۔ آخری چیز~$\TMstroke$ دیاں لڑیوں دا اک سلسلہ
  اے جو~$M_e$ دے نشاناں دے کوڈ عدداں، یعنی
  مصنوعی “نشاناں”، دی نمائندگی کردا اے؛ لڑیوں نوں
  ~$\TMblank$ وکھ کردا اے۔

  ہن ایہ مشین اوہ عمل نقل کردی اے جیڑا $M_e$
  انپٹ~$n$ اُتے شروع ہو کے کردِی؛ اوہ ایہ قدم چکدی اے:
  \begin{enumerate}
    \item “سر دی موجودہ جگہ” دا عدد~$k$ لبھو
    (شروع وچ ایہ~$1$ اے)؛
    \item “پٹی” دے $k$-ویں ٹکڑے ول جا کے اُتھوں دا
    “نشان” ویکھو؛
    \item\ollabel{find-inst}%
    موجودہ “حالت” تے “نشان” نال میل کھاندی ہدایت لبھو؛
    \item “پٹی” دے $k$-ویں ٹکڑے ول واپس جا کے اوتھے
    دے “نشان” دی تھاں اُس نشان دا کوڈ عدد لکھو
    جیڑا $M_e$ لکھدی؛
    \item سر نوں اُس تھاں لے جاؤ جتھے موجودہ
    “حالت” درج اے، تے اوتھے دا عدد نویں حالت
    دے عدد نال بدل دو؛
    \item جتھے “پٹی دی جگہ” درج اے اُتھے جا کے
    اک~$\TMstroke$ مٹا دو یا اک~$\TMstroke$ ودھا دو
    (ترتیب نال، جے ہدایت کھبّے یا سجّے ہلّن لئی
    آکھے)؛
    \item ایہ کم دوبارہ کرو۔\footnote{ایتھے اسی کجھ باریک
    اوکھائیاں دا پورا حساب نہیں دِتا۔ مثلاً “سر دی
    موجودہ جگہ” دا حساب رکھن والی گنتی ودھے تاں
    $U$~نوں ہور جگہ چاہیدی ہو سکدی اے؛ اوس صورت
    وچ اوہنوں ساری مصنوعی “پٹی” سجّے پاسے
    سرکاؤنی پَوے گی۔}
  \end{enumerate}
  جے $M_e$ انپٹ~$n$ اُتے شروع ہو کے کدی نہیں
  رُکدی، تاں $U$ وی نہیں رُکدی؛ ایس لئی اوہدا
  آؤٹ پٹ تعریف شدہ نہیں۔

  جے قدم~\olref{find-inst} وچ پتا لگے کہ~$M_e$ دے
  بیان وچ موجودہ “حالت”/“نشان” جوڑی لئی کوئی ہدایت
  نہیں، تاں $M_e$ رُک جاندی۔ ایسا ہووے تاں $U$
  اپنی پٹی دا مصنوعی “پٹی” دے کھبے پاسے والا حصہ
  مٹا دیندی اے۔ تین~$\TMstroke$ نشاناں دی ہر لڑی
  لئی، جیڑی اک~$\TMstroke$ دی نمائندگی کردی اے
  (ایہ نشان $M_e$ دی پٹی اُتے اے)، اوہ اپنی پٹی دے
  کھبے سِرے اُتے اک $\TMstroke$ لکھدی تے نال نال مصنوعی
  “پٹی” نوں مٹاندی جاندی اے۔ ایہ کم مکّے تاں
  $U$ دی پٹی اُتے~$\TMstroke$ نشاناں دی
  لمبائی~$m$ والی اکّو لڑی رہ جاندی اے۔
  
  جے $U$ نوں مصنوعی “پٹی” اُتے تین~$\TMstroke$
  نشاناں دی لڑی توں وکھ کجھ ملے، تاں اوہ فوراً
  رُک جاندی اے۔ ایس صورت وچ $U$ دی پٹی اُتے
  ~$\TMstroke$ نشاناں دی اکّو لڑی نہیں ہوندی؛
  ایس لئی اوہدا آؤٹ پٹ قدرتی عدد نہیں، یعنی
  ایس صورت وچ $f(e,n)$ تعریف شدہ نہیں۔
\end{proof}

\end{document}
